\bigskip

\begin{enumerate}
\item  Le plan $\mathcal{P}_1$ a pour équation $5x+2y+4z=17$,  un vecteur normal à $\mathcal{P}_1$ est donc $\vect{n_1}=\begin{pmatrix} 5\\2\\4 \end{pmatrix}$. 

De même le plan $\mathcal{P}_2$ a pour équation $10x+14y+3z=19$,  un vecteur normal à $\mathcal{P}_2$ est  $\vect{n_2}=\begin{pmatrix} 10\\14\\3 \end{pmatrix}$. 

Il n'existe pas de réel $\lambda$ tel que $\vect{n_2}=\lambda \vect{n_1} \quad \left(\dfrac{10}{5}\neq\dfrac{14}{2}\right)$  donc les vecteurs $\vect{n_1}$ et $\vect{n_2}$ ne sont pas colinéaires, d'où les plans $\mathcal{P}_1$ et $\mathcal{P}_2$ ne sont pas parallèles.

\item Les plans $\mathcal{P}_1$ et $\mathcal{P}_2$ ne sont pas parallèles,  ils sont donc sécants.

Vérifions que la droite $\mathcal{D}$ est l'intersection des plans  $\mathcal{P}_1$ et $\mathcal{P}_2$.

Soit $t$ un réel quelconque et soit $M_t$ le point de paramètre $t$ de la droite $\mathcal{D}$.

$5x_{M_t}+2y_{M_t}+4z_{M_t}=5 (1+2t)+2\times (-t) +4 (3 - 2t)= 5+10t-2t+12-8t=17$ donc $M_t$ est aussi un point du plan $\mathcal{P}_1$.

$10x_{M_t}+ 14y_{M_t}+ 3z_{M_t}=10 (1 + 2t)+14\times (- t) +3 (3 - 2t)= 10+20t-14t+9-6t=19$ donc $M_t$ est aussi un point du plan $\mathcal{P}_2$.

$M_t$ appartient donc l'intersection de $\mathcal{P}_1$ et $\mathcal{P}_2$,  la droite $\mathcal{D}$ est donc bien  l'intersection des plans  $\mathcal{P}_1$ et $\mathcal{P}_2$.

\item
\begin{enumerate}
\item   $5x_{\mathrm{A}}+2y_{\mathrm{A}}+4z_{\mathrm{A}}=5 \times 1 +2\times (-1) +4 \times (-1)= 5-2-4=-1 \neq 17$ donc le points A n'appartient pas au plan $\mathcal{P}_1$.
	\item  A appartient à la droite $\mathcal{D}$ si et seulement si il existe un réel $t$ tel que : 
	
$\left\{\begin{array}{@{}l} x_{\mathrm{A}}=1+2t\\y_{\mathrm{A}}=-t\\z_{\mathrm{A}}=3-2t\end{array}\right. 	
\iff \left\{\begin{array}{@{}l} 1=1+2t\\-1=-t\\-1=3-2t\end{array}\right. 	  
\iff \left\{\begin{array}{@{}l} 0=2t\\1=t\\2t=4\end{array}\right. 	  
\iff \left\{\begin{array}{@{}l} t=0\\t=1\\t=2\end{array}\right. $

Ce système n'admet pas de solution donc le point A n'appartient donc pas à la droite $\mathcal{D}$.
\end{enumerate}

\item
	\begin{enumerate}
		\item  Soit $t$ un réel :

$f(t)=\mathrm{A}M^2= {(x_M - x_\mathrm{A})^2+(y_M - y_\mathrm{A})^2+(z_M - z_\mathrm{A})^2} ={ (1+2t-1)^2+ (-t+1)^2+(3-2t+1)^2    }$

$f(t) ={ (2t)^2+ (-t+1)^2+(4-2t)^2  } 
={ 4t^2+ t^2-2t+1+16-16t+4t^2    }= 9t^2-18t+17$ 

	\item La distance A$M$ est minimale si et seulement si la distance A$M^2$ est minimale c'est à dire $f(t)$ est minimale.

	$f$ est un polynôme du second degré avec un coefficient dominant (9)  positif.  $f$ admet donc un minimum pour $t=\dfrac{-(-18)}{2\times 9} = \dfrac{18}{18}=1$.
	
	Le point de $\mathcal{D}$ qui correspond au paramètre $t=1$ est le point $M$ de coordonnées $(1+2\times 1~;~-1~;~3-2\times 1)$ c'est à dire $M(3~;~-1~;~1)$, ce qui est ce qu'il fallait démontrer.

	\end{enumerate}

\item Un vecteur directeur de la droite (AH) est le vecteur $\vect{\mathrm{AH}} \begin{pmatrix} x_\mathrm{H}-x_\mathrm{A}\\y_\mathrm{H}-y_\mathrm{A}\\z_\mathrm{H}-z_\mathrm{A} \end{pmatrix} = \begin{pmatrix} 3-1\\-1+1\\1+1\end{pmatrix}= \begin{pmatrix} 2\\0\\2\end{pmatrix}$

D'après la représentation paramétrique de la droite $\mathcal{D}$,  un vecteur directeur de la droite $\mathcal{D}$ est $\vect{u}\begin{pmatrix}2\\-1\\-2 \end{pmatrix}$ .

Comme le repère est orthonormé,  on peut calculer les produits scalaires à l'aide des coordonnées :

$ \vect{u}\cdot \vect{\mathrm{AH}} = 2\times 2 - 1 \times 0 + 2 \times (-2) = 4 + 0 - 4 = 0$ : $\vect{u}$ est orthogonal à $ \vect{\mathrm{AH}}$.

Les droites (AH) et $\mathcal{D}$ sont donc orthogonales.

De plus, d'après la question précédente, le point H appartient à la droite $\mathcal{D}$, il appartient aussi à la droite (AH), ces deux droites sont donc sécantes au point H, elles sont donc perpendiculaires.
\end{enumerate}

\bigskip
